\section{delta 的几种选择}

脚本里把用于更新的量统一记成 $\delta$，然后尝试了几种常见做法：

\begin{itemize}
    \item \textbf{rewards}：直接用 raw rewards；
    \item \textbf{centered rewards}：减去 group mean；
    \item \textbf{normalized rewards}：再除以标准差；
    \item \textbf{max rewards}：只保留组内最大 reward。
\end{itemize}

\begin{knowledgebox}{最常见的是 centered rewards}
在同一个 prompt 的多个 response 里，减掉 group mean 是最自然的 baseline 近似。它对应“相对同组平均表现的好坏”。
\end{knowledgebox}

\begin{table}[H]
\centering
\begin{tabular}{@{}lll@{}}
\toprule
\textbf{模式} & \textbf{公式} & \textbf{教学含义} \\
\midrule
\texttt{rewards} & $\delta_i = r_i$ & 直接用绝对 reward，最朴素但方差最大 \\
\texttt{centered\_rewards} & $\delta_i = r_i - \bar r$ & 组内相对比较，最像 baseline \\
\texttt{normalized\_rewards} & $\delta_i = (r_i - \bar r)/(\sigma + \epsilon)$ & 再压缩尺度，训练更平滑 \\
\texttt{max\_rewards} & 只保留组内最大值 & 强调 winner-take-all 的更新方向 \\
\bottomrule
\end{tabular}
\caption{脚本中四种 $\delta$ 选择的差别：本质上都是在决定“谁应该被强化”}
\end{table}


